
\documentclass[12pt,french]{article}
\usepackage[utf8]{inputenc}
\usepackage[left=5in,right=5in,top=5in,bottom=5in]{geometry}
\usepackage{bm}
\usepackage{amsmath}
\usepackage{amssymb}
\usepackage{indentfirst}
\usepackage[hyperpageref]{backref} % back references biblio
\usepackage{tocbibind}
\usepackage[round,sort&compress]{natbib}
\setcitestyle{aysep={}}
\usepackage{amsfonts}
\usepackage{enumerate}
\usepackage{babel}
\usepackage{caption}
\usepackage{supertabular}
\usepackage{tabularx}
\usepackage{float}
\usepackage{dsfont}
\usepackage{fancyvrb}
\usepackage{verbatim}
\usepackage[hyphens]{url}
\usepackage{hyperref}
\usepackage[shortlabels]{enumitem}
\usepackage{setspace}
\usepackage{comment}
\usepackage{subcaption}
\usepackage{graphicx}
\usepackage{tikz}
\usetikzlibrary{shapes,backgrounds,positioning}
\usepackage{gensymb}
\usepackage{eurosym}
\usepackage{textcomp}
\usepackage{color,soul}

\usepackage{multicol}
\usepackage{changepage}
\usepackage{enumitem}


\usepackage{tabulary}
\usepackage{tabularx}
\usepackage{booktabs}
\usepackage{fullpage}
\usepackage{morefloats}
% \usepackage[utf8]{inputenc}
% \usepackage{bm}
% \usepackage{indentfirst}
% \usepackage{tocbibind}
% \usepackage{enumerate}
\usepackage{makecell}
\usepackage{multirow}
\usepackage{ulem}
\usepackage{lscape}
\usepackage{pdflscape}
\usepackage{longtable}
\usepackage{rotating}
\usepackage{fancyhdr}
\usepackage{tocloft}
\usepackage{multibib}
\usepackage{titletoc}
\usepackage[export]{adjustbox}
\usepackage[anythingbreaks]{breakurl} % for links
%\usepackage[nomarkers,figuresonly]{endfloat} % Figures at the end
\hypersetup{
  colorlinks=true, % breaklinks=true,
  urlcolor=blue, % color of external links
  linkcolor=blue,  % color of toc, list of figs etc.
  citecolor=violet,   % color of links to bibliography
}

% Commande pour distinguer visuellement les réponses aux commentaires
\newcommand{\reponse}[1]{\medskip\par\noindent\textbf{Réponse~:} #1\par\medskip}

\title{Préférences budgétaires des Français~: résultats d'une enquête représentative ~\\ ~\\ \textbf{Réponses aux rapporteurs}}

\date{\today}
\begin{document}

\maketitle

\noindent Je remercie les deux rapporteurs pour leurs lectures attentives et leurs commentaires constructifs. Je reproduis ci-dessous chacun de leurs commentaires (en italique) et y réponds point par point. Les modifications correspondantes dans le manuscrit sont signalées le cas échéant.

\bigskip

%============================================================
\paragraph*{Rapporteur \#1}
%============================================================

\textit{Cet article présente une analyse originale de préférences budgétaires des Français à partir d'une enquête. L'analyse est descriptive mais des observations intéressantes sont faites et la visualisation des résultats est innovante. Je trouve l'article intéressant et je suis très favorable à la démarche. Je détaille ci-dessous ce que je considère comme les principales faiblesses de la version actuelle de l'article.}

Je tiens à remercier le rapporteur pour son rapport favorable, diligent et constructif. Je réponds ci-dessous à chacun de ses commentaires, en indiquant les modifications apportées dans le manuscrit.

\subparagraph*{Points essentiels}

\bigskip
\noindent\textbf{1.} \textit{Je pense qu'il serait utile de fournir davantage de détails sur les questions posées dans l'enquête ainsi que sur sa méthodologie, idéalement dans une annexe dédiée. Cela permettrait au lecteur de mieux évaluer la qualité des données et de comprendre précisément le dispositif empirique. Par ailleurs, je recommande de présenter plus explicitement la méthodologie empirique propre à chacun des exercices dans le corps du texte. Même une courte sous-section dédiée avant chaque analyse faciliterait la lecture, et rendrait l'enchaînement des résultats plus fluide et plus accessible.}

\reponse{Je remercie le rapporteur pour cette suggestion. Pour compléter la présentation du questionnaire en Section 1, j'ai ajouté une annexe dédiée à la présentation de l'enquête, qui reproduit les captures d'écran des deux questions principales (\texttt{budget} et \texttt{programme}) telles qu'elles ont été soumises aux répondants. Le lecteur peut ainsi vérifier la formulation exacte des questions, ainsi que le compteur d'économies affiché en continu lors de l'évaluation des mesures budgétaires. Par ailleurs, j'ai veillé à ce que le début de chaque section détaille la méthodologie empirique propre à l'exercice correspondant.}



\bigskip
\noindent\textbf{2.} \textit{Je pense qu'il est important de clarifier le cadre conceptuel d'économie politique qui sous-tend l'analyse. À plusieurs reprises, le lecteur a du mal à identifier le mécanisme théorique auquel les résultats sont censés se rapporter. Dans certaines parties du texte, l'analyse semble renvoyer à un modèle de formation de coalitions (par exemple, lorsque l'auteur évoque une alliance entre «~la gauche~» et des «~nativistes frugaux~»). Ailleurs, notamment dans la Section~5, l'interprétation paraît davantage relever d'un cadre d'électeur médian. D'autres passages semblent plutôt faire référence à des processus de négociation législative. Cette ambiguïté est particulièrement visible lorsque les réponses individuelles aux différentes mesures sont interprétées. Le lecteur peut avoir l'impression que l'exercice correspond à une situation où chaque proposition serait soumise séparément à référendum. Pourtant, l'interprétation retenue ne correspond pas entièrement à ce cadre (voir NBP~15 p.19). L'approche proposée est volontairement a-théorique, ce qui est un choix tout à fait défendable. En revanche, il me semble qu'un ancrage plus explicite dans la littérature serait utile. En particulier, il serait intéressant de discuter les travaux portant sur la formation de plateformes politiques multidimensionnelles, ainsi que les modèles spatiaux du vote, afin de mieux situer la contribution de l'article et d'éclairer l'interprétation des résultats empiriques.}

Je remercie le rapporteur pour cette remarque. J'assume le caractère volontairement a-théorique de la démarche, mais j'ai ajouté dans l'introduction un paragraphe qui ancre explicitement l'interprétation des résultats dans la littérature sur les modèles spatiaux du vote à programmes multidimensionnels~:
\begin{quote}
\textit{«~Si le cadre d'analyse de cette étude est a-théorique, certains résultats peuvent s'interpréter à l'aide des modèles spatiaux du vote avec programmes multidimensionnels. En particulier, l'insatisfaction de nombreux Français envers l'ensemble des partis peut refléter le théorème du chaos de McKelvey qui énonce que bien souvent, aucun vainqueur de Condorcet ne se dégage  (Davis et al., 1970; Downs, 1957; Enelow \& Hinich, 1984; McKelvey, 1976).~»}
\end{quote}
Ce cadre offre une clé de lecture unifiée~: l'absence de vainqueur de Condorcet éclaire aussi bien la difficulté à faire émerger une majorité stable autour d'un paquet unique de mesures que la coexistence, selon les dimensions considérées, de coalitions différentes. Il ne s'agit toutefois pas de tester un mécanisme théorique précis, mais de situer la contribution empirique par rapport à cette littérature.



\bigskip
\noindent\textbf{3.} \textit{J'ai une difficulté avec l'exercice de formation d'une majorité autour de paquets de mesures. L'auteur indique~: «~Pour comprendre quel budget pourrait être porté par une coalition majoritaire, il ne s'agit pas seulement de juxtaposer les 97~Mds de mesures convenables pour différentes majorités, mais de trouver un paquet de mesures acceptées conjointement par la même majorité~». Cette distinction est importante et constitue le cœur de l'exercice proposé. Toutefois, même si l'on accepte cet exercice comme une analyse fondée sur des préférences révélées, il demeure problématique que l'enquête ne fournisse aucune observation sur un choix effectif portant sur un ensemble de mesures supposées être acceptées conjointement. Ma compréhension de l'enquête est que les répondants évaluent les mesures individuellement, mais ne sont jamais confrontés à des arbitrages entre différents paquets de politiques publiques. Dès lors, le passage des préférences exprimées sur des mesures prises isolément à l'acceptation d'un paquet cohérent de réformes repose sur une hypothèse forte, qui mériterait d'être davantage discutée et justifiée. En l'état, il me semble que les conclusions de cet exercice sont plus fortes que ce que les données permettent d'établir.}

Je concède que le cadre d'analyse employé permet d'identifier le paquet de mesures idéal pour un répondant et les paquets majoritaires au sein de groupes d'électeurs, mais ne permet pas d'étudier l'arbitrage qu'effectueraient les électeurs entre certains programmes donnés (ceux des partis politiques). En outre, l'exercice ne tient pas compte de l'importance relative de différentes mesures aux yeux des répondants, qui pourrait affecter les préférences collectives entre différents programmes. J'ai ajouté la mise en garde suivante en Section~5 pour éviter que les conclusions soient surinterprétées~: \textit{Notons que le plus grand paquet accepté par une majorité ne coïncide pas nécessairement avec le paquet préféré des Français, puisque l'exercice ne tient pas compte de l'importance relative de différentes mesures aux yeux des répondants, qui pourrait affecter les préférences collectives entre différents programmes. }

\bigskip
\noindent\textbf{4.} \textit{L'auteur propose, dans plusieurs sections de l'article, des labels pour caractériser différents groupes de répondants (par exemple, «~sociaux-nativistes~», «~frugaux~», «~progressistes~», etc.) à partir d'un ensemble de réponses à quelques questions. Je pense que cette terminologie peut prêter à confusion, en particulier pour des lecteurs familiers d'autres classifications ou d'autres usages de ces concepts dans la littérature. Je suggère donc d'être plus explicite sur la manière dont ces labels sont construits et sur ce qu'ils sont censés représenter. Lorsque cela est possible, il serait également utile de les mettre en perspective avec la littérature existante en économie politique et en science politique, en utilisant une terminologie établie ou, à défaut, en expliquant clairement en quoi les catégories proposées s'en distinguent. Cela faciliterait l'interprétation des résultats et éviterait d'attribuer aux groupes des caractéristiques qui dépassent ce que les données permettent réellement d'identifier.}

Je remercie le rapporteur pour cette remarque, qui rejoint mon souci d'éviter les sur-interprétations. J'ai ajouté, à l'endroit où les profils sont nommés (Section~6), une note de bas de page qui précise explicitement le statut et la construction de ces étiquettes~: «~\textit{Les noms donnés aux profils (}conservateurs\textit{, }progressistes\textit{, }frugaux\textit{, puis plus loin \textit}libéraux-nativistes\textit{, }sociaux-nativistes\textit{ et }sociaux\textit{) sont des étiquettes choisies a posteriori pour résumer les attitudes qui caractérisent chaque profil issu du partitionnement~: soutien majoritaire ou minoritaire aux grandes catégories de mesures (taxation des plus aisés, État-providence, mesures nativistes, impôts indifférenciés). Elles ne renvoient ni à des catégories partisanes ni à une typologie établie de la littérature. En particulier, contrairement à la tripartition (en blocs }social-écologique\textit{, }libéral-progressiste\textit{ et }national-patriote\textit{) de Cagé \& Piketty (2023) fondée sur les programmes des partis, les profils étudiés ici émergent des préférences uniquement programmatiques et peuvent regrouper des électeurs d'appartenances partisanes différentes -- ainsi les {sociaux-nativistes} réunissent-ils des électeurs à la fois favorables aux mesures nativistes et à la redistribution, combinaison que ne capture aucune offre partisane. Si l'analyse de Cagé \& Piketty (2023) a l'avantage de reposer sur des préférences révélées, elle identifie les préférences des électeurs au programme de leur parti et est incapable d'observer un décalage entre les deux, tel que le soutien majoritaire des électeurs centristes à des mesures redistributives ou nativistes, à rebours des }programmes\textit{ centristes.~»}



\bigskip
\noindent\textbf{5.} \textit{La discussion Section~7 consacrée à la validité des mesures, à leurs effets attendus ou encore à leur conformité constitutionnelle ne me semble pas être directement centrée sur l'objet de l'étude. Il n'est pas toujours clair quel est son statut dans l'argumentation. L'auteur formule-t-il ces remarques parce qu'il considère que les répondants ont une interprétation des mesures différente de celle retenue par le CAE~? S'agit-il d'une mise en garde sur la manière dont les résultats doivent être interprétés~? Ou bien ces développements relèvent-ils d'une discussion plus générale sur les politiques publiques considérées~? En l'état, cette section donne parfois davantage l'impression d'un commentaire sur la situation politique que d'une discussion académique des résultats. Par exemple, des affirmations telles que «~le bloc central aura du mal à gagner sur une promesse de sueur et de larmes~» ou encore «~Électoralement, le RN ferait sans doute une erreur s'il ne proposait pas de programme budgétaire crédible.~» apparaissent relativement éloignées de l'analyse empirique présentée dans l'article. Je suggère donc de recentrer cette section sur le cœur de la contribution. Si ces développements sont destinés à éclairer l'interprétation des résultats empiriques, il serait utile de rendre ce lien plus explicite. Dans le cas contraire, leur portée pourrait être réduite afin de maintenir le focus sur les questions de recherche et les enseignements directement étayés par les données.}

La discussion finale relève il est vrai de l'analyse politique des enseignements de l'étude, et va au-delà de l'interprétation académique des résultats empiriques. Je souhaiterais conserver cette analyse car elle fait partie intégrante de la motivation de l'étude. Toutefois, j'ai recentré cette analyse en retirant ou reformulant les passages mentionnés, qui pouvaient être perçus comme des conseils destinés aux différents blocs politiques. J'ai également spécifié au début de la discussion que : «~\textit{
Suite à l'analyse empirique des résultats, il est possible de proposer quelques enseignements utiles au débat public.}~» 

\subparagraph*{Points non essentiels}

\bigskip
\noindent\textbf{6.} \textit{L'auteur mentionne en passant dans une note de bas de page qu'il a réalisé des entretiens~: «~Outre les données de l'enquête, ce paragraphe est informé par des entretiens individuels menés auprès d'un échantillon représentatif de 100~Français.~». Aucune information additionnelle n'est donnée sur la méthode utilisée ou les résultats, ce qui est frustrant pour le lecteur.}

Ces entretiens vont faire l'objet de plusieurs publications à venir. Comme ils sont ici utilisés de façon relativement anecdotique, il serait inopportun d'en détailler la méthodologie et les résultats. J'ai cependant étendu la note pour indiquer~: «~\textit{Outre les données de l'enquête, ce paragraphe est informé par des entretiens individuels de 50 minutes menés auprès d'un échantillon représentatif de 100 Français sur leurs attitudes politiques, et notamment par leurs réponses à la question~:} Pouvez-vous énumérer les principaux partis politiques français et dire ce que vous pensez de chacun d'entre eux ?.~»

\bigskip
\noindent\textbf{7.} \textit{L'abstract m'a semblé trop long et gagnerait à être plus concis et plus informatif. Il s'agit en partie d'une question de préférence rédactionnelle, raison pour laquelle je classe cette remarque parmi les points non essentiels. Néanmoins, dans sa version actuelle, l'abstract ne met pas suffisamment en avant la question de recherche, la méthodologie, les principaux résultats et la contribution de l'article. Une version plus resserrée permettrait au lecteur de comprendre plus rapidement l'apport du papier.}



\bigskip
\noindent\textbf{8.} \textit{Pourquoi le choix du vote aux élections européennes de 2024, et non les législatives suivantes, plus proches de la question étudiée~?}



\bigskip
\noindent\textbf{9.} \textit{Sur le tableau~1 je m'attends également à un commentaire sur l'extrême droite.}



\bigskip
\noindent\textbf{10.} \textit{Coquilles~: p.1 NBP~1 «~Antoine~» $\Rightarrow$ Antonin~; p.2 «~à gauche.2~» NBP avant le point.}

\reponse{[À compléter. Corrections effectuées dans le manuscrit.]}

\bigskip

%============================================================
\paragraph*{Rapporteur \#2}
%============================================================

\textit{Cet article présente les résultats d'une enquête d'opinion concernant les préférences budgétaires des Français.}

\subparagraph*{Principaux commentaires}

\bigskip
\noindent\textbf{2.0.} \textit{Cet article est très intéressant. Il s'intéresse à une question importante. La méthodologie est claire. L'article est bien rédigé. Les résultats sont informatifs pour la politique française. Cependant, il ne s'agit pas, selon moi, d'un article de recherche. Sa contribution scientifique n'est pas apparente. Cela n'enlève rien à la qualité de l'article ou à son intérêt.}



\bigskip
\noindent\textbf{2.} \textit{Les calculs utilisant les écarts-moyens ou les écarts (distances) entre répondants sont questionnables car il s'agit de réponses qualitatives ordinales. Tout est traité comme si chaque mouvement d'une «~graduation~» avait la même valeur. On peut évidemment faire cette hypothèse, mais elle est assez problématique, car le passage du côté négatif d'une évaluation au côté positif (par exemple, de «~moins favorable~» à «~plus favorable~») a la même valeur que le passage de la situation neutre à une évaluation positive (par exemple, de «~ne changerait rien~» à «~beaucoup plus favorable~»). Le même commentaire s'applique au passage de «~inacceptable~» à «~supportable~» et au passage de «~supportable~» à «~convenable~». Peut-être serait-il souhaitable de montrer que des regroupements en deux catégories ne modifient pas drastiquement l'analyse. Dans tous les cas, cette hypothèse de travail doit, à minima, être soulignée (elle l'est trop brièvement en page~12).}



\bigskip
\noindent\textbf{2.2.} \textit{Les abstentionnistes sont identifiés au début de l'article, mais ils ne sont pas du tout analysés. Or, l'abstention est différente aux élections européennes (utilisées pour identifier ce groupe) et aux élections nationales. L'absence de cette prise en compte est dommageable, compte tenu de l'objectif principal de l'article.}



\subparagraph*{Autres commentaires}

\bigskip
\noindent\textbf{2.3.} \textit{Les électeurs d'extrême droite constituent le groupe majoritaire après exclusion des abstentionnistes. Il est donc mécanique que leurs attitudes soient celles qui sont les plus proches de celles de l'ensemble des électeurs (page~12).}



\bigskip
\noindent\textbf{2.4.} \textit{Je n'ai pas du tout compris l'explication des 50~\%, 32~\% et 68~\% en page~7.}



\bigskip
\noindent\textbf{2.5.} \textit{La note de bas de page 14 (page~18) n'est pas claire.}



\bigskip
\noindent\textbf{2.6.} \textit{La section~6 est très intéressante~!}



\clearpage
\renewcommand{\url}[1]{\href{#1}{Link}}
% \bibliographystyle{plainnaturl_clean}
% \bibliography{global_tax_attitudes}

\end{document}
